\section{Logarithmic curves, line bundles and stable maps}
\subsection{Curves and line bundles}
We
refer the reader to \cite{ogus} for an extensive introduction to logarithmic structures and in particular for all the basic definitions which we do not recall in what follows.
\subsubsection{Logarithmic curves}
Let $(S,M_S)$ be a fine and saturated logarithmic scheme. A log curve over $S$ is a proper, integral and logarithmically smooth morphism $\pi\colon C\to S$ with connected, reduced one-dimensional (geometric) fibers.
Kato provided the following local characterization.
\begin{theo}[\cite{KatoF}]
For every geometric point $p\in C$ with $\pi(p)=s$, there exists an \'etale local neighbourhood of $p$ in $C$ with a strict \'etale morphism to
\begin{description}\itemsep=0pt
 \item[\emph{smooth point:}] $\mathbb A^1_S$ with log structure pulled back form the base, i.e., $\overline{M}_{C,p}=\overline M_{S,s}$;
 \item[\emph{marking:}] $\mathbb A^1_S$ with log structure generated by the $0_S$-section and $\pi^*M_S$, i.e., $\overline{M}_{C,p}=\overline M_{S,s}\oplus \mathbb N$;
 \item[\emph{node:}] $\operatorname{Spec}(\mathcal O_S[x,y]/xy-t)$ for some $t\in\mathcal O_S$, with log structure induced by the multiplication map $\mathbb A^2_S\to\mathbb A^1_S$ and $t\colon S\to \mathbb A^1$. In particular, $\overline{M}_{C,p}=\overline M_{S,s}\oplus \mathbb N e_x\oplus \mathbb N e_y\slash e_x+e_y=\delta$ for~${\delta\in \overline{M}_{S,s}}$ the so called \emph{smoothing parameter}.
\end{description}
\end{theo}

\subsubsection{Tropical curves}

To a family of logarithmic curves, we can associate a family of tropical curves.
Consider first $C\to (\operatorname{Spec}(k),Q)$ a logarithmic curve over a log point.

The \emph{tropicalization} $\Gamma$ is defined as the generalised cone complex obtained as the following colimit:
\[
\varinjlim_{\eta\rightsquigarrow\xi} \operatorname{Hom}\bigl(\overline{M}_{C,p},\mathbb R_{\geq 0}\bigr),
\]
where $\eta\rightsquigarrow\xi$ denote specialization maps inducing surjective morphisms $\overline{M}_{C,\xi}\to\overline{M}_{C,\eta}$. From the local description of the log structure on log smooth curves recalled above, it follows that $\Gamma$ comes equipped with a morphism of cone complexes
$\Gamma\to\sigma_Q $,
for $\sigma_Q=\operatorname{Hom}(Q,\RR_{\geqslant 0})$ the cone dual to the monoid $Q$.

Alternatively, $\Gamma$ can be thought as the dual graph of $C$ (having a vertex for each irreducible component, a leg for each marking and an edge $e\in\Gamma$ for each node) together with a metrization of the bounded edges with values in the base monoid $Q$. This means that we have an assignment~${l\colon E_b(\Gamma)\to Q}$ defined by associating to each bounded edge its smoothing parameter~${\delta_e\in Q}$.

\begin{remarkk}
Notice that for each map $Q\to\mathbb R_{\geq 0}$ we obtain a tropical curve in the classical sense. Since $Q\to\mathbb R_{\geq 0}$ correspond to a point of the cone $\sigma_Q$ dual to $Q$ (in the usual sense of toric geometry), we can think of a tropical curve metrized in $Q$ as a family of usual tropical curves on the dual cone.
\end{remarkk}
\begin{remarkk}\label{rem:tropfamilies}
The tropicalization makes sense over more general base log schemes (see also \cite[Definition~2.3.3.3]{{molchowiselog}}). Let $(C,M_C)\to(S,M_S)$ a log smooth curve;
for each $s\in S$, $\Gamma_s$ consists of the dual graph of $C_s$ metrized in $\overline{M}_{S,s}$.
 For each specialization $\eta\rightsquigarrow s$, we have a compatible diagram
\[
\begin{tikzcd}
\Gamma_s\ar[r,"c"]\ar[d,"\ell"] &\Gamma_{\eta}\ar[d,"\ell"]\\
\overline{M}_{S,s}\ar[r] & \overline{M}_{S,\eta},
\end{tikzcd}
\]
where $\overline{M}_{S,s}\to \overline{M}_{S,\eta}$ is a localization morphism followed by sharpification (quotient the invertibles) and $c$ is an edge contraction; more precisely, $c$ contracts the edges of $\Gamma_s$ whose length becomes zero in $ \overline{M}_{S,\eta}$.

\end{remarkk}
\subsubsection{Piecewise linear functions}
Unravelling the definition of global sections of the characteristic sheaf, we get a convenient description of $H^0\bigl(C,\overline{M}_C^{\rm gp}\bigr)$ in terms of continuous functions on the tropicalization $\Gamma$ which are linear with integral slope along the edges.
Let us first assume that $C$ is a log curve over a~${S=(\operatorname{Spec}(k),Q)}$.

Since $\overline{M}_C^{\rm gp}$ is a constructible sheaf, a
section is determined by giving $m_x\in\overline{M}_{C,x}^{\rm gp}$ at a generic point $x$ of each stratum,
such that the values $m_x$ are compatible with specialization, i.e., if $\xi$ generalize to $\eta$ then the image of $m_{\xi}$ under the surjection $\overline{M}_{C,\xi}^{\rm gp}\to \overline{M}_{C,\eta}^{\rm gp}$ is $m_{\eta}$.

Looking at Kato's classification, there are three type of strata on log smooth curves: the open strata corresponding to irreducible components of the curves excluding markings, i.e., vertices in the tropicalization; the closed strata corresponding to markings; the closed strata corresponding to nodes. Therefore, a section $\alpha\in H^0\bigl(C,\overline{M}_C^{\rm gp}\bigr)$ is the data of
\begin{itemize}\itemsep=0pt
 \item for each vertex $V$ of $\Gamma$ a value $\alpha(V)\in Q$;
 \item for each leg and edge an integer $w_e\in\mathbb Z$ (the slope along edge or leg in $\Gamma$);
 \item a compatibility condition for the values $\alpha(V)$, $\alpha(W)$ of vertices incident to the same edge~$e$: $\alpha(V)-\alpha(W)=\delta_e w_e$.
\end{itemize}
Indeed, suppose that $x_e$ is a node
of $C$ which generalizes to the generic points $\eta_V$ and $\eta_W$; then we have a map
$\overline{M}_{C,x_e}^{\rm gp}\to Q\times Q$
that is injective and gives an identification
\[\overline{M}_{C,x_e}^{\rm gp}=\{(a,b)\in Q\times Q \mid a-b\in\delta_e\mathbb Z\}.\]

We denote by $\operatorname{CL}(\Gamma,Q)$ the set of functions satisfying the above conditions, called \textit{cone-wise linear}.\footnote{We reserve the name piece-wise linear for those function which become linear along the edges after a subdivision.} The previous discussion can be summarized in the following equality (see also \cite[Section~3.3.3]{marcuswiselog} and references therein)
$H^0\bigl(C,\overline{M}_C^{\rm gp}\bigr)=\operatorname{CL}(\Gamma,Q)$.

 \begin{remarkk}
 Using the description of tropicalization for a family $(C,M_C)\xrightarrow{\pi}(S,M_S)$ of log curves given in Remark~\ref{rem:tropfamilies} and the above discussion, we can think of a section $ \alpha\in H^0\bigl(S,\pi_*\overline{M}_C^{\rm gp}\bigr)$ as a collection \smash{$\alpha_s\in \operatorname{CL}\bigl(\Gamma_s,\overline{M}_{S,s}^{\rm gp}\bigr)$} compatible under edge contraction, i.e., generalization.
 \end{remarkk}



\subsubsection{Line bundles}
Let $C\xrightarrow{\pi} S$ be a log smooth curve. We have a short exact sequence
\begin{equation}\label{eq:logseq}
0\to\mathcal O_C^{\times}\to M_C^{\rm gp}\to\overline{M}_C^{\rm gp}\to 0,\end{equation}
from which we obtain a long exact sequence of sheaves on $S$
\begin{align*}
 \cdots\to \pi_*M_C^{\rm gp}\to\pi_*\overline{M}_C^{\rm gp}\to
 R^1\pi_*\mathcal O_C^{\times}\to
 R^1\pi_*M_C^{\rm gp}\to R^1\pi_*\overline{M}_C^{\rm gp}\to\cdots.
\end{align*}
In particular, if the underlying scheme of $S$ is a point, so that $\pi_*\overline{M}^{\rm gp}_C=H^0\bigl(C,\overline{M}^{\rm gp}_C\bigr)$ and $R^1\pi_*\O_C^\times=H^1\bigl(C,\O_C^\times\bigr)=\Pic(C)$, we get a map
\smash{$
 H^0\bigl(C,\overline{M}_C^{\rm gp}\bigr)\to \Pic^0(C)$}, $
 \alpha \mapsto\mathcal O_C(-\alpha)$.

\begin{remarkk}
 The sign comes from observing that when $\alpha\in H^0\bigl(C,\overline{M}_C\bigr)$, the $\mathcal O_C^{\times}$-torsors of lifts of $\alpha$ to a section of $M_C$ comes equipped with a map to $\mathcal O_C$, coming from the structure morphism~${\epsilon \colon M_C\to\mathcal O_C}$ of the log structure. Up to changing $\alpha$ by $-\alpha$, we now deal with $\O_C(\alpha)$ instead.
\end{remarkk}

The restriction of the line bundle $\mathcal O_C(\alpha)$ to the irreducible components $C_V$ of $C$ can be explicitly described \cite[Proposition~3.3.3]{marcuswiselog}.
\begin{prop}
 Let $C\xrightarrow{\pi} S=(\operatorname{Spec}(k),Q)$ be a log curve and $\alpha\in\operatorname{CL}(\Gamma, Q)$, where $\Gamma$ denotes the tropicalization of $C$. Then for any vertex $V$,
 \[\mathcal O_C(\alpha)\rvert_{C_V}=\pi^*\mathcal O_S(\alpha(V))\otimes\mathcal O_{C_V}\biggl(\sum_{e\vdash V} s_e(\alpha) q_e\biggr),
 \]
where the sum is taken over all the edges and legs $e$ incident to $V$ and $s_e(\alpha)$ is the slope of~$\alpha$ along~$e$, oriented away from~$V$, where $q_e$ denotes the preimage of the node associated to the edge~$e$ on the considered component.
\end{prop}
\begin{remarkk}
 The result is more generally true if $C\to S$ is a logarithmic curve with constant degeneration.
\end{remarkk}

\subsection{Logarithmic line bundles and trivializations}
\begin{defi}[\cite{molchowiselog}]
 Let $C\to S$ be a log smooth curve. A \emph{logarithmic line bundle} $P$ on $C$ is a $M_C^{\rm gp}$-torsor such that its restriction $P\rvert_{C_s}$ to the geometric fibers has \emph{bounded monodromy}.
\end{defi}
We refer the reader to \cite{molchowiselog} for all explanations about the meaning and the necessity of the bounded monodromy. Here we only need that it is a condition on the $\overline{M}_C^{\rm gp}$-torsor $\overline{P}$ associated to $P$, and that it is automatically satisfied when $\overline{P}$ is isomorphic to the trivial $\overline{M}_C^{\rm gp}$-torsor.

In particular, we see from the long exact sequence
\[H^0\bigl(C,\overline{M}_C^{\rm gp}\bigr)\to\Pic(C)\to H^1\bigl(C,M_C^{\rm gp}\bigr)\to H^1\bigl(C,\overline{M}_C^{\rm gp}\bigr)\to\cdots\]
that for any $L\in \Pic(C)$ the associated $M_C^{\rm gp}$-torsor \smash{$L^{\times}\otimes_{\mathcal O_C^{\times}}M_C^{\rm gp}$} is in fact a logarithmic line bundle, where we denote by $L^{\times}$ the $\mathbb{G}_m$-torsor obtained from the line bundle removing its zero section.

\begin{defi}[{\cite[Definition~4.5.1]{marcuswiselog}}]
 Let $C\xrightarrow{\pi} S$ be a log smooth curve and $L$ a line bundle on $C$. A \emph{logarithmic trivialization} of $L$ is a section of \smash{$L^{\times}\otimes_{\mathcal O_C^{\times}}M_C^{\rm gp}$}.
\end{defi}
The definition is a modification of \cite[Definition~4.5.1]{marcuswiselog}, where the authors define logarithmic trivializations relative to the base. The difference comes from the fact that \cite{marcuswiselog} are interested in moduli of logarithmic maps to rubber targets, while we study the rigid case.
Similarly, the following discussion is parallel to
\cite[Proposition~4.5.3]{marcuswiselog}.

From \eqref{eq:logseq}, we see that the logarithmic line bundle \smash{$L^{\times}\otimes_{\mathcal O_C^{\times}}M_C^{\rm gp}$} can be thought as a \smash{$\mathcal O_C^{\times}$}-torsor over the trivial $\overline{M}_C^{\rm gp}$-torsor.
In particular, there is an exact sequence
\[0\to H^0\bigl(C,L^{\times}\otimes_{\mathcal O_C^{\times}}M_C^{\rm gp}\bigr)\to H^0\bigl(C, \overline{M}_C^{\rm gp}\bigr)\to\Pic(C),
 \]
where the second map is defined by
$\alpha\mapsto L(-\alpha)=L\otimes\O(-\alpha)$.
In other words, a logarithmic trivialization of $L$ is a trivialization of the line bundle $L(-\alpha)$ for some cone-wise linear function~$\alpha$.

\subsection[Logarithmic stable maps to P\^{}1-bundles and their degenerations]{Logarithmic stable maps to $\boldsymbol{\mathbb P^1}$-bundles and their degenerations}

This subsection follows closely \cite[Section~5]{marcuswiselog}, with two small differences. First, \cite{marcuswiselog} considers maps to rubber target, while we deal with the rigid case. Second, they study maps to $\mathbb{P}^1$ or to degeneration of $\mathbb P^1$ to a chain, while we consider maps to $\mathbb{P}^1$-bundles $Y=\mathbb P_X(\mathcal O_X\oplus L)$ and degenerations of the latter to certain chains $\Y_0=\bigl(Y_1,D_1^\pm\bigr)\cup \bigl(Y_2,D_2^\pm\bigr)\cup\dots\cup\bigl(Y_N,D_N^\pm\bigr)$ of~$\mathbb{P}^1$-bundles over $X$ obtained by gluing $Y_i$ to $Y_{i+1}$ along $D_i^+\cong D_{i+1}^-\cong X$.

For the benefit of the reader, we state the results in the form most convenient for us and sketch how to adapt the proofs of \cite[Section~5]{marcuswiselog} to our situation.

Let $X$ be a smooth projective variety and $L$ a line bundle on $X$.
For a fixed $\beta\in H_2(X,\mathbb Z)$ and a vector of integers $\bfw=(w_1,\dots,w_n)$ such that $\sum_{i=1}^n w_i=c_1(L)\cdot\beta$, we consider the moduli space
$\M := \M_{g,m}\bigl(Y|D^{\pm},\beta,\bfw\bigr)$
of logarithmic stable maps to $Y=\mathbb P_X(\mathcal O_X\oplus L)$ endowed with the divisorial log structure $D^++D^-=\mathbb P_X(\mathcal O_X) +\mathbb P_X(L)$ and with prescribed contact order $w_i$ at $q_i$ along $D^++D^-$ \cite{abramovich2014stable,gross2013logarithmic}.
Positive $w_i$ encode the tangencies with $D^+$ and the negative $w_i$ the tangencies with $D^-$.

This space compactifies the moduli space parametrizing maps from marked curves to $X$,
\[f\colon\ \bigl(C,p_1,\dots, p_m,q_1,\dots,q_n\bigr)\to X,\]
together with an isomorphism
\[\lambda\colon\ f^* L\biggl(-\sum a_i q_i\biggr)\cong\mathcal O_C.\]

Using the language of logarithmic trivializations introduced in the previous subsection, this description can be extended over the boundary.

\begin{defi}\label{def:logstablemaps}
 We define $\M(X,L)$ to be the stack over $\operatorname{LogSch}^{fs}$ whose objects are
 \begin{enumerate}\itemsep=0pt
 \item [(1)] A log smooth curve $C\to S$ with $n$ marked points with logarithmic structure and $m$ schematic marked points.
 \item [(2)] A stable map $f\colon (C,p_i,q_i)\to X$.
 \item [(3)] A logarithmic trivialization \smash{$\lambda\in H^0\bigl(C,f^*L^{\times}\otimes_{\mathcal O_C^{\times}}M_C^{\rm gp}\bigr)$}, such that the induced $\alpha\in H^0\bigl(C,\allowbreak\overline{M}_C^{\rm gp}\bigr)$ is locally on $C$ comparable with $0$ and the slope of $\alpha$ along the $i$-th marking is $w_i$.
 \end{enumerate}

We recall that a section \smash{$\alpha\in H^0\bigl(C,\overline{M}^{\rm gp}_C\bigr)$} is said comparable to $0$ if at each point $x\in C$ we have that either $\alpha(x)\geq 0$ or $\alpha(x)\leq 0$ in the partial order on $\overline{M}^{\rm gp}_{C,x}$ defined by the monoid $\overline{M}_{C,x}$. We refer the reader to \cite[Section~2.5]{RSPW2} for further explanations on the geometric meaning of this condition.
\end{defi}

\begin{prop}\label{prop:logstabelmaps}
The moduli space $\M(X,L)$ is equivalent as a stack over logarithmic schemes to the moduli space of logarithmic stable maps $\M$ of Gross--Siebert, Abramovich--Chen with target~${(Y=\PP_X(\mathcal O_X\oplus L),\mathcal M_Y)}$ and fixed tangencies $w_i$ along the markings $q_1,\dots,q_n$ {\rm\cite{abramovich2014stable,gross2013logarithmic}}. Here~$\mathcal M_Y$ is the divisorial log structure defined by the fiberwise $0$ and $\infty$ sections.
\end{prop}
\begin{proof}
The proof essentially consists in unravelling the definition of logarithmic stable map to a $\mathbb P^1$-bundle. For the case of log stable maps to toric varieties, the analogous result is proved, for example, in \cite[Proposition~2.5.1.1]{RSPW2}.

Let $\mathrm{Tot}(L)\xrightarrow{\pi} X$ be the total space of $L$ endowed with the zero section logarithmic structure~$M_L$. First, we argue that a logarithmic stable map
\smash{$C\xrightarrow{F} \mathrm{Tot}(L)$} is the same as a logarithmic trivialization $\lambda$ such that the associated $\alpha$ is a section of the log structure, i.e., $\alpha\in H^0\bigl(C,\bar{M}_C\bigr)$.
Indeed, by definition, a log map is a map $\underline{F}$ of the underlying schemes, namely $f\colon C\to X$ plus a section of $f^*L$, together with a morphism $\underline{F}^{-1}M_L\to M_C$ of the logarithmic structures. Using the Borne--Vistoli description of fine and saturated log structures \cite{BorneVistoli}, the latter is the same as a~morphism of the characteristic sheaves $\underline{F}^{-1}\bar{M}_L\to \bar{M}_C$ such that the maps to $\mathfrak{D}\mathrm{iv}_C$ commute. The first map comes from
$
 \bar{M}_L\cong\mathbb N_X\to \mathfrak{D}\mathrm{iv}_{\mathrm{Tot}(L)}$, $
1\to \mathcal O_L(X)\cong \pi^* L
$
pulling back along $\underline{F}$. It follows that a lift of the morphism of characteristic sheaves to a~morphism of the logarithmic structures correspond to an isomorphism $f^*L\cong^{\lambda}\mathcal O_C(\alpha)$ for $\alpha$ the section of $\bar{M}_C$ induced by the map on characteristic sheaves.

At this point we argue parallel to \cite[Section~2.5.1]{RSPW2}.

The $\mathbb P^1$-bundle $\mathbb P_X(\mathcal O_X\oplus L)$ with its fiberwise toric log structure can be constructed as the quotient of the rank 2 vector bundle $\mathcal O_X\oplus L$ minus the $0$ section for the fiberwise $\mathbb G_m$-action.

Then any logarithmic map $C\to \mathbb P_X(\mathcal O_X\oplus L)$ locally lifts to $\O_X\oplus L\setminus 0_X$.

By the discussion above, this lift can be represented by $a$, $b$ logarithmic trivializations of~$\mathcal O_C$ and $f^*L$ such that the associated $\bar{M}_C^{\rm gp}$ section $\bar{a}$ and $\bar{b}$ are in $\bar{M}_C$. Notice that the ratio~${\lambda=ab^{-1}}$ is a well defined $\mathbb G_m$-invariant section of the torsor \smash{$f^*L^{\times}\otimes_{\mathcal O_C^{\times}}M_C^{\rm gp}$}, namely a logarithmic trivialization of $f^*L$.

Notice however that not all the section $\alpha\in \rm{H}^0\bigl(C,\bar{M}_{C}^{\rm gp}\bigr)$ can arise this way: indeed, since $(a,b)$ must avoid the zero section, this means that locally around each point $x\in C$ either $\bar{a}_x\in \bar{M}_{C,x} $ or $\bar{b}_x\in\bar{M}_{C,x}$ is zero (i.e., lifts to an isomorphism of trivial $\mathbb G_m$-torsors). For $\alpha=\bar{a}-\bar{b}$, this translates precisely to being locally comparable to zero.
\end{proof}

\begin{remarkk}
If in Definition~\ref{def:logstablemaps}, we drop the request on the local comparability with zero of the section $\alpha\in H^0\bigl(C,\bar{M}_C^{\rm gp}\bigr)$, we would obtain the stack over $\operatorname{LogSch}^{fs}$ parametrizing log maps to the $\mathbb G_{log}$ bundle $\rm{Tot}(L)\times_{\mathbb G_m}\mathbb G_{log}$ associated with $L$. Similar moduli stack have been considered various times in the literature \cite{marcuswiselog,ranganathan2024logarithmic,RSPW2,ranganathan2020rational}.\footnote{We thank an anonymous referee whose comments suggested us to add this remark.}
\end{remarkk}

\begin{remarkk}
The definition and the characterization above can easily be adapted to the case of moduli of logarithmic stable maps to expansions \cite{kim2010logarithmic,li2001stable,maulik2023logarithmic}. Indeed, it is sufficient to add the condition that the values of $\alpha(V)$ are totally ordered; similar discussions appear in \cite{marcuswiselog,ranganathan2024logarithmic}.

In the classical language of \cite{li2001stable}, maps in the boundary of $\M$ which differ by the rubber action are identified. It is explained in detail in \cite{CN} how to see this from the logarithmic prospective.
\end{remarkk}

We are also interested in considering log maps to degenerations, i.e.,
\begin{itemize}\itemsep=0pt
\item Logarithmic maps to \smash{$\Y_t\xrightarrow{p}\mathbb A^1_t$} certain one parameter degenerations of $\mathbb P_{\X_{t\neq 0}}(\O\oplus\L)$ to
\[\Y_0=\bigl(Y_1,D_1^\pm\bigr)\cup\bigl(Y_2,D_2^\pm\bigr)\cup\dots\cup \bigl(Y_N,D_N^\pm\bigr),
\]
where $\bigl(Y_i,D^{\pm}_i\bigr)$ are $\mathbb P^1$-bundles over $\X_0$ glued along their infinity and zero sections;
$\mathbb A^1_t$ is endowed with its toric log structure and $\Y_t$ with the divisorial log structure defined by~${\Y_0+\mathcal{D}^+ +\mathcal{D}^-}$, where $\mathcal{D}^\pm$ are the $0$ and $\infty$-section of the family, and making $p$ log smooth.
\item Logarithmic maps to the log smooth scheme $\Y_0\to(\operatorname{Spec} k,\mathbb N)$ with the log structure pulled back from the family $\Y_t$
\end{itemize}

We recall the following definition.

\begin{defi}[{\cite[Definition~5.2.1]{marcuswiselog}}]
 A \emph{divided tropical line} over a logarithmic scheme $(S,M_S)$ is defined as follows. Let $P$ be an $M_S^{\rm gp}$-torsor and $\bar{P}$ the associated $\bar{M}_S^{\rm gp}$ torsors; fix $\gamma_0\leq\gamma_1\leq\cdots\leq\gamma_N$ a non empty collection of sections of $\bar{P}$ defined locally on $S$. Then $\bar{P}_{\gamma}\subset\bar{P}$ is defined to be the subfunctor whose local sections are comparable with all the $\gamma_i$.
\end{defi}
It is proved in \cite[Proposition~5.2.4]{marcuswiselog} that $P_{\gamma}:=\bar{P}_{\gamma}\times_{\bar{P}} P$ is a 2-marked family of semi-stable genus zero curve.



We then consider the following variation: fix $(\mathcal X, M_{\X})\xrightarrow{p} (S,M_S)$ with $X\to S$ smooth projective and $M_{\X}=p^*M_S$.
The two cases of interest listed above correspond to the case where~$(S,M_S)$ is $\mathbb A^1$ with its toric log structure or a standard log point.

Let $P\in\operatorname{LogPic}_{\X\slash S}$ and
$\bar{P}$ the induced $\bar{M}^{\rm gp}_{\mathcal{X}}$-torsor. Let $\gamma_0\leq\gamma_1\leq\cdots\leq\gamma_N$ be a non empty collection of sections of \smash{$\bar{P}$}. Notice that where the logarithmic structure of $\X$ is trivial $P$ is a line bundle in the usual sense and all the $\gamma_i$ coincide and are necessarily zero.
\begin{lem}\label{lem:tropicaline}
In the notation above, let $\bar{P}_{\gamma}\subseteq \bar{P}$ be the sub-funtcor of sections locally comparable with all the $\gamma_i$ and $P_{\gamma}=P\times_{\bar{P}}\bar{P}_{\gamma}$. Then $P_{\gamma}\to\X\to S$ is an expanded degeneration of $\Y=\mathbb P_{\X}(\L\oplus\O)$ for $\L$ the line bundle on $\X$ of lifts of $\gamma_0$ to a section of $P$.
\end{lem}
\begin{proof}
Since we assumed that the collection of sections is not empty, $\bar{P}$ admits a trivialization~${\gamma_0\colon\X\to\bar{P}}$ which we think as the zero of $\bar{M}_{\X}$ group. Notice that if the log structure on~$S$ and thus on $\X$ is generically trivial, then all the $\gamma_i$ coincide and are actually zero on the locus of $\X$ with trivial logarithmic structure. Then there exist a line bundle $\L_0$ representing the log line bundle $P$, namely the $\O_{\X}^\times$-torsor of lifts of $\gamma_0$ to a section of the log line bundle. We already proved in Proposition~\ref{prop:logstabelmaps} that the subfunctor of \smash{$P=L_0^\times\otimes_{\O_{\X}^\times}M_{\X}^{\rm gp}$} whose local sections are~comparable with $\gamma_0$ is the $\mathbb{P}^1$-bundle $\Y=\mathbb P_{\X}(\L_0\oplus\O)$ endowed with its boundary logarithmic structure.
Once a trivialization is fixed, we have that
\[
\gamma_i-\gamma_0=\delta_i\in H^0\bigl(\bar{M}_{\X}^{\rm gp}\bigr);
\] these are maps \smash{$\operatorname{Trop}(\X)=\operatorname{Trop}(S)\xrightarrow{\delta_i}\operatorname{Trop}(\Y)$}. Subdividing along the image, we obtained the desired degeneration.

 Alternatively, the statement can be proved following the steps of \cite[Proposition~5.2.4]{marcuswiselog}.
\end{proof}

Let $\mathcal Y_t\to\mathbb A^1$ a semi-stable degeneration coming form a subdivided tropical line; for example, this is the case if $\Y_t$ is constructed starting from $\mathbb P_{\X_t}(\O_{\X_t}\oplus\L_t)$ by successively blowing up the zero and the infinity section on the central fiber.
Denote by $\M\bigl(\Y_t\slash\mathbb A^1\bigr)$ the moduli space of logarithmic stable maps to $\Y_t\slash\mathbb A^1$ from families of log smooth curves with fixed genus $g$, number of markings $m$, contact order $(w_1,\dots, w_n)$ along $\D_t^-+\D_t^+$ and curve class $\beta$. Then we have the following.

\begin{coro}\label{prop:logmapsdegeneration}
 The moduli space $\M\bigl(\Y_t\slash\mathbb A^1\bigr)$ parametrizes the following data:
 \begin{enumerate}\itemsep=0pt
 \item[$(1)$] A diagram of logarithmic maps
 \[
 \begin{tikzcd}
 C\ar[r,"f"]\ar[d] & \X_t\ar[d]\\
 S\ar[r] &\mathbb A^1
 \end{tikzcd}
 \]
 for $C\to S$ a family of log smooth curves.
 \item[$(2)$] A section \smash{$\lambda\in H^0\bigl(C, f^* \mathcal L^{\times}\otimes_{\mathcal O^{\times}_{C}}M_{C}^{{\rm gp}}\bigr)$} such that the induces $\alpha\in H^0\bigl(C,\bar{M}_C^{\rm gp}\bigr)$ is locally comparable with $f^*\gamma_i$.
 \end{enumerate}
\end{coro}
\begin{remarkk}
 Parallel to before, if we want to instead consider logarithmic maps to expanded degenerations \cite{kim2010logarithmic,li2001stable,maulik2023logarithmic}, it suffices to further impose that for each geometric point $s$ the values of $\alpha(V)$ for $V\in\Gamma_s$ are totally ordered.
\end{remarkk}

\section{Correlated virtual class}

\subsection{Recollection on Albanese varieties}

\subsubsection{Complex setting}
For $X$ a proper, smooth variety over the complex numbers the Albanese variety was originally defined via transcendental methods using path integrals of closed holomorphic 1-forms. Let~$H_1(X,\mathbb Z)$ denote the torsion free part of the first homology group. There is an inclusion via the integration on paths:
\begin{align*}
 H_1(X,\mathbb Z)\longrightarrow H^0(X,\Omega_X)^*,\qquad
 \gamma\longmapsto \left(\omega\mapsto\int_{\gamma}\omega\right),
\end{align*}
where $ H^0(X,\Omega_X)^*$ denotes the dual to the vector space of holomorphic $1$-forms.
The Albanese variety of $X$ is the complex abelian variety obtained as the following quotient:
\[\Alb(X):=H^0(X,\Omega_X)^*\slash H_1(X,\mathbb Z). \]
We list below some results and properties of the Albanese variety and refer the reader to \cite[Section~11.11]{birkenhake} or \cite{mumford1974abelian} and references therein for all details and proofs.
\begin{enumerate}\itemsep=0pt
 \item[(1)] When $X$ is an Abelian variety, then $\Alb(X)\cong X$.
 \item[(2)] Let $x_0\in X$ a point, then there is an algebraic map, called the \emph{Albanese map}
 \begin{align*}
 a_{X}\colon \ X\longrightarrow \Alb(X),\qquad
 x\longmapsto\left(\omega\mapsto\int_{x_0}^x\omega\right) \mod H_1(X,\mathbb Z),
 \end{align*}
 sending $x_0$ to the identity element. This map is unique up to translation in the sense that choosing a different $x_0$ amounts to compose $a_X$ with a translation.
 \item[(3)] The Albanese map satisfies the following universal property:
 let $\varphi\colon X\to A$ be a map to an abelian variety, then there exists a unique homomorphism of abelian varieties $\widetilde{\varphi}\colon \Alb(X)\to A $ such that $\widetilde{\varphi}(0)=\varphi(x_0)$ and $\widetilde{\varphi}\circ a_{X}=\varphi$.
 \item[(4)] Let $X\xrightarrow{f} Y$ a morphism of algebraic varieties, then the map to the Albanese varieties are functorial, i.e., there exists a homomorphism of abelian varieties $f_*$ making the following diagram commute:
 \[
 \begin{tikzcd}
 X \arrow[r,"f"] \arrow[d,"a_X"']& Y \arrow[d,"a_Y"] \\
 \Alb(X) \arrow[r,"f_*"] & \Alb(Y),
 \end{tikzcd}
\]
 where $a_Y$ maps $f(x_0)$ to the identity.
 \item[(5)] The Albanese variety $\Alb(X)$ is the dual abelian variety of the Picard group
 \[\Pic^0(X):=H^1(X,\mathcal O_X)\slash H^1(X,\mathbb Z), \]
 i.e.,
 $\Alb(X)=\Pic^0\bigl(\Pic^0(X)\bigr)$.
 Identifying $\Pic^0(X)$ with the connected component of the identity of the group of line bundles, the functorial homomorphisms of Albanese variety~$f_*$ is the dual of the pullback map.
\end{enumerate}

\begin{remarkk}
 Notice that in the complex setting, $H^1(X,\O_X)$ can be seen as sheaf cohomology or alternatively as the Dolbeault cohomology group of $(0,1)$-forms.
\end{remarkk}



\subsubsection{Algebraic definition}
The interpretation of the Albanese variety as the dual abelian variety of \smash{$\operatorname{Pic}^0_{X\slash S}$}
allows a more algebraic definition.
 Indeed, when $\X\xrightarrow{g} S$ be a smooth, projective, morphism with connected geometric fibers over a scheme $S$, Grothendieck proved that there exists a smooth $S$-scheme, locally of finite presentation, $\Pic_{\X\slash S}$ representing the relative Picard functor and an abelian scheme (in particular, smooth and proper) $\Pic^0_{\X\slash S}$ over $S$ whose fibers $\Pic^0_{\X_s\slash k(s)}$ are the connected components of the identity in $\Pic_{\X_s\slash k(s)}$. See, for example, \cite[Section~9]{fantechi2006fundamental}, or \cite[Chapter~8]{bosch2012neron} and references therein.
Then by \cite[Theorem~3.3]{grothendieck1961technique} (see also \cite{milneAV,mochizuki2012topics}), we can consider the dual abelian scheme
\[\Alb^0_{\X\slash S}:=\Pic^0_{\Pic^0_{\X/S}\slash S}=\bigl(\Pic^0_{\X\slash S}\bigr)^\vee.\]
If $g\colon \X\to S$ has a section $\sigma_S\colon S\to\X$, then we get an Albanese map to $\Alb^0_{\X\slash S}$ constructed in the following way.

Let $\mathcal P$ be the \emph{unique} Poincar\'e line bundle on $\X\times\Pic^0_{\X\slash S}$ which trivializes along $\sigma_S$ in the sense of \cite[Section~4]{bosch2012neron}.

Then we define the Albanese morphism
\[
 a_{\X\slash S}\colon \ \X \longrightarrow \Alb^0_{\X\slash S},\qquad
 \big(U\xrightarrow{f} \X\big) \longmapsto (f\times\mathrm{id})^*\mathcal P.
 \]
Given \smash{$\X\xrightarrow{f} \Y$} a morphism of smooth, proper, geometrically connected $S$-varieties and $\sigma_S\colon S\allowbreak\to \X$ a section, it follows from the universal property of the Albanese scheme that there exists a~unique homomorphism $f_*$ of $S$-abelian schemes making the following diagram commute:
\[
\begin{tikzcd}
\X\ar[r,"f"]\ar[d,"a_{\X\slash S}"'] & \Y\ar[d,"a_{\Y\slash S}"]\\
\Alb^0_{\X\slash S}\ar[r,"f_*"] & \Alb^0_{\Y\slash S},
\end{tikzcd}
\]
where $a_{\X\slash S}$ and $a_{\Y\slash S}$ sends the sections $\sigma_S\colon S\to \X$ respectively $f\circ \sigma_S \colon S\to \Y$ to the identity element. Alternatively, $f_*$ can be described as the dual (in the category of abelian $S$-schemes, see \cite[Section~I.8]{milneAV}) of the pullback
$f^*\colon \Pic^0_{\Y\slash S}\to \Pic^0_{\X\slash S}$.


\subsubsection{Self-duality for smooth curves} As explained, for example, in \cite[Section~9.5.26]{fantechi2006fundamental}, when $\C\to S$ is a family of smooth, projective, geometrically connected genus $g>0$
curves over a Noetherian base, there exists an \emph{self-duality isomorphism}
$\Alb^0_{\C\slash S}\to \Pic^0_{\C\slash S}$.

 If we have a section $\sigma_0\colon S\to \C$, then the self-duality has a very explicit description using the Abel--Jacobi map $A_{\sigma_0}\colon\C\to\Pic^0_{\C/S}$ given by
\[
 A_{\mathcal \sigma_0}\colon \ \C \longrightarrow \Pic^0_{\C\slash S},\qquad
 (T\to\C) \longmapsto \mathcal O_{\C_T}(p_T -\sigma_0\rvert_T),
\]
where $p_T\colon T\to\C\times_S T$ is the induced section of $\C_T=\C\times_S T$.
The self-duality isomorphism is simply given by the pullback along the Abel--Jacobi map. The isomorphism does not depend on the choice, nor on the existence of the section, see, for example, \cite{esteves2002autoduality} for a proof.

Notice that $A_{\sigma_0}$ coincides with the Albanese map defined using the Poincar\'e line bundle on~${\X\times\Pic^0_{\X\slash S}}$ rigidified along the section $\sigma_0$.

\subsection{Refinement of the moduli spaces}

Let $\M=\M_{g,m}\bigl(Y|D^\pm,\beta,\bfw\bigr)$ be the moduli space of logarithmic stable maps to $Y=\mathbb P_X^1(\O\oplus L)$ considered in the previous section and let us now assume that $L\in\Pic^0(X)$: we have $c_1(L)=0$.


\subsubsection[Morphism to Alb(X)]{Morphism to $\boldsymbol{\operatorname{Alb}(X)}$}
Using the tangency orders $\bfw$, we consider the map
\[ a_\bfw\colon \ (x_i)\in X^n\longmapsto\sum_{i=1}^n w_i a_X(x_i)\in\Alb(X). \]
As $\sum w_i=0$, this map does not depend on the choice of a base point in $X$ and is uniquely defined. Indeed, for $z_0$, $z_0'$ two choices of base point and $a_X$, $a'_X$ the corresponding morphisms, we have that $a'_X= t_{a'_X(z_0)}\circ a_X$ and thus
\[
 \sum_{i=1}^n w_i a'_X(x_i)= \sum_{i=1}^n w_i (a_X(x_i) -a'_X(z_0))= \sum_{i=1}^n w_i a_X(x_i).
\]
Now, consider the morphism from the moduli space to the Albanese variety of $X$ by composing with the evaluation map:
\[ \kappa\colon\ \M\xrightarrow{\ev} X^n\xrightarrow{a_\bfw} \Alb(X)\]
defined by
\[%\label{eq:refinement morphism}
\bigl(\C\xrightarrow{f} X,(p_j)_{j=1}^m,(q_i)_{i=1}^n,\O_{\C}(\alpha)\cong^{\lambda} f^*L\bigr)\mapsto (x_i:=f\circ q_i)_{i=1}^n\mapsto \sum_{i=1}^n w_i a_X(x_i).
\]

Recall that $(q_i)_{i=1}^n$ are marking with logarithmic structure while $(p_j)_{j=1}^m$ are standard schematical points on $\C$.

\begin{lem}\label{lem:mapphibeta}
Let $\beta\in H_2(X,\ZZ)$ be the homology class of a complex curve. The bilinear map
\[ \phi \colon\ \alpha\otimes\omega\in H^{0,1}(X)\otimes H^{1,0}(X)\longmapsto \int_\beta \alpha\wedge \omega\in\CC \]
induces a morphism
\[ \varphi_\beta \colon\ \Pic^0(X)\longrightarrow \Alb(X). \]
\end{lem}
\begin{proof}
We use the isomorphisms $H^{0,1}(X)=H^1(X,\O_X)$ and $H^{1,0}(X)=H^0(X,\Omega_X)$ as well as the Hodge decomposition $H^2(X,\CC)=H^{0,1}(X)\oplus H^{1,0}(X)$. Through the above isomorphisms, the map $H^1(X,\ZZ)\to H^1(X,\O_X)$ yielding $\Pic^0(X)$ is actually the composition of $H^1(X,\ZZ)\to H^1(X,\CC)$ followed by the projection onto $H^{0,1}(X)$.
The bilinear map $\phi$ induces a~linear map~${\phi^*\colon H^{0,1}(X)\to H^{1,0}(X)^*}$. To prove that it descends to a map $\Pic^0(X)\to\Alb(X)$, we need to show that
\[ \phi^*\bigl(H^1(X,\ZZ)\bigr)\subset H_1(X,\ZZ)\subset H^0(X,\Omega_X)^*. \]
Let $\omega\in H^{1,0}(X)$ and $\gamma\in H^1(X,\ZZ)$. We write $\gamma\otimes\CC=\gamma^{0,1}+\gamma^{1,0}$ for its Hodge decomposition in $H^2(X,\CC)$. We have
\begin{align*}
\phi^*(\gamma)(\omega) = {}& \int_\beta \gamma^{0,1}\wedge\omega \qquad\text{(by definition)}\\
={} & \int_\beta \bigl(\gamma^{0,1}+\gamma^{1,0}\bigr)\wedge\omega
= \beta\cap (\gamma\cup\omega)
= (\beta\cap\gamma)\cap\omega,
\end{align*}
where the second equality follows from the fact that \smash{$\int_\beta \gamma^{1,0}\wedge\omega=0$}, since $\beta$ is the class of a~complex curve. The last line ensures that $\phi^*(\gamma)(\omega)$ is a period of $\omega$ since $\beta\cap\gamma\in H_1(X,\ZZ)$. Therefore, as desired, $H^1(X,\O_X)\to H^0(X,\Omega_X)^*$ descends to a map
\[ \varphi_\beta \colon\ H^1(X,\O_X)/H^1(X,\ZZ)\longrightarrow H^0(X,\Omega_X)^*/H_1(X,\ZZ). \tag*{\qed}
\]\renewcommand{\qed}{}
\end{proof}

\begin{remarkk}
Giving a non-degenerate positive bilinear map $H^1(X,\O_X)\otimes H^0(X,\Omega_X)\to\CC$ amounts to providing a polarization, i.e., an ample line bundle $\L$, on the Abelian variety $\Pic^0(X)$, see, for example,~\cite{mumford1974abelian}. In an Abelian variety, a polarization then induces a map
\[ A_\L\colon\ L\in\Pic^0(X)\longmapsto t_L^*\L\otimes \L^{-1}\in\Pic^0\bigl(\Pic^0(X)\bigr)=\Alb(X). \]
Assume $X$ is projective with hyperplane class $h\in H^2(X,\ZZ)$ giving a K\"ahler form on $X$. Then we have a polarization given by the non-degenerate positive bilinear form
\[ \phi\colon\ \alpha\otimes\omega\longmapsto \int_X h^{n-1}\wedge\alpha\wedge\omega. \]
Assume the class $\beta$ is obtained by intersecting sufficiently many hyperplane sections: $\beta$ is Poincar\'e dual to $h^{n-1}$. Then, the previously defined morphism~$\varphi_\beta$ is actually the morphism~$A_\L$ provided by the above polarization on $\Pic^0(X)$.
\end{remarkk}
We claim that the morphism $\kappa$ is in fact constant to $\varphi_\beta(L)$. To prove the latter, we start with the case of smooth curve, and then use it to deal with the case of nodal curves.

\begin{prop}\label{prop-image-of-L-is-constant}
The morphism $\kappa\colon\M\rightarrow \operatorname {Alb}^0(X)$ defined above is constant equal to $\varphi_\beta(L)$.
\end{prop}

\begin{proof}
Let us first look at a geometric point $s=\operatorname{Spec}(\mathbb C)\to\M$ of the moduli space such that~the source curve $C$ is smooth. We claim that in such case
$\kappa(s)=f_*f^*L$,
where~${\O_C \bigl(\sum_{i=1}^n w_i q_i\bigr)\cong f^*L}$ (by definition of the moduli functor) and
$f_*\colon\Pic^0(C)\to\operatorname{Alb}^0(X)$ is the morphism obtained dualizing $f^*\colon\Pic^0(X)\to\Pic^0(C)$ via the self-duality isomorphism of~$\Pic^0(C)$.

To see that, choose $p_0$ a point in $C$ and $z_0=f(p_0)$; then the functoriality of the Albanese morphism for compatible choices of base point gives
\[f_*(f^*L)=f_*\biggl(\O_C\biggl(\sum_{i=1}^n w_i q_i\biggr)\biggr)=f_*\biggl(\sum_{i=1}^n w_i A_{p_0}(q_i)\biggr)= \sum_{i=1}^n w_i a_X(x_i)\]
with the last equality following from the commutativity.

Since we are working over $\mathbb C$, we may use that for $X$ (resp.\ $C$), we have
 \[ \Pic^0(X)=H^1(X,\O_X)/H^1(X,\ZZ) \qquad \text{and}\qquad \Alb(X)=H^0(X,\Omega_X)^*/H_1(X,\ZZ). \]
Through Hodge theory, $H^1(X,\O_X)$ is isomorphic to the Dolbeault cohomology group $H^{0,1}(X)$. In particular, by functoriality, the pullback map $f^*\colon\Pic^0(X)\to\Pic^0(C)$ is in fact induced by~${f^*\colon H^1(X,\O_X)\to H^1(C,\O_C)}$, which is the pullback at the level of $(0,1)$-forms. Using that $H^0(X,\Omega_X)=H^{1,0}(X)$, the composition $f_* f^*$ is induced by the $\CC$-linear map
 \[ \varphi_C \colon\ H^{0,1}(X) \xrightarrow{f^*} H^{0,1}(C)\simeq H^{1,0}(C)^*\xrightarrow{f_*}H^{1,0}(X)^*, \]
where the second arrow is the dual to the pullback map for holomorphic forms. The middle isomorphism is provided by the Poincar\'e duality. Finally, if $\alpha\in H^{0,1}(X)$ is a $(0,1)$-form on $X$ and $\omega\in H^{1,0}(X)$ an holomorphic $1$-form, we have
 \[ \varphi_C(\alpha)(\omega) = \int_C f^*\alpha\wedge f^*\omega = \int_C f^*(\alpha\wedge\omega) = \int_{f(C)}\alpha\wedge\omega= \int_\beta \alpha\wedge\omega, \]
where the penultimate equality is push-pull formula.

Now, assume that $C$ is nodal; we write $\nu\colon\bigsqcup C_V\to C$ for the normalization map and $f_V\colon C_V\to X$ for the restriction of $f\circ\nu$ to the component $C_V$ of the normalization. Since each $C_V$ is now smooth, the discussion above tells us that the morphism
\[
f_{V\ast}f_V^*\colon\ \Pic^0(X)\to\Pic^0(C_V)\to\operatorname{Alb}^0(X)
\]
is equal to $\varphi_{\beta_V}$, where $\beta_V=f_{V\ast}[C_V]$. Furthermore, since $\sum \beta_V=\beta$, we have
\begin{align*}
\varphi_{\beta}=\sum \varphi_{\beta_V}\colon\ \Pic^0(X) \to\prod_V\Pic^0(C_V)\to\operatorname{Alb}^0(X),\qquad
 L \to \nu^*f^*L\to \sum_V f_{V\ast}(f_V^*L).
\end{align*}

To conclude the proof, we only need to argue that
\[ \sum_V f_{V\ast}(f_V^*L)=\kappa(s).\]

When $C$ is singular, $f^*L\cong^{\lambda}\O_C(\alpha)$ for $\alpha\in H^0\bigl(C, \bar{M}_C^{\rm gp}\bigr)$ such that
\[\O_C(\alpha)\rvert_{C_V}=\O_{C_V}\biggl(\sum_{e\vdash V}s_{e,V}(\alpha)q_e\biggr),\]
satisfying the balancing condition $\sum_{e\vdash V}s_{e,V}(\alpha)=0$. Here $s_{e,V}(\alpha)$ is the outgoing slope of $\alpha$ at~$V$. Notice that $\alpha$ induces an orientation on the edges of $\Gamma$. We denote with $w_e$ the slope along an edge with this induced orientation. Then if $V$, $W$ are the two vertices adjacent to a~bounded edge $e$, we will have $s_{e,V}(\alpha)=-s_{e,W}(\alpha)$ with one of the two being $w_e$.

 We know from the analysis for the smooth case that
 \[f_{V\ast}(f_V^*L)=\sum_{e\vdash V}s_{e,V}(\alpha) a_X(f_V\circ q_i).\]

 From the observation above and the fact that $f_V(q_e)=f_W(q_e)$ since the maps are induced by the maps on the nodal curves, it follows that when we consider $\sum_Vf_{V\ast}(f_V^*L)$ the contributions coming from the nodes cancel out and we obtain the required identity.
\end{proof}

\begin{expl}
 Assume $X=E$ is an elliptic curve. We have $\Alb(E)\cong\Pic^0(E)\cong E$. For the homology class $\beta=a[E]$, the associated map $\varphi_\beta$ is just the multiplication by $a$.
\end{expl}

\begin{remarkk}
For a family \smash{$C\slash S\xrightarrow{F} X\times S$} with non-smooth family of source curves, it is not obvious how to define
\[F_*\colon\ \Pic^{[0]}_{C\slash S}\to\operatorname{Alb}^0(X)\times S.\]
Indeed, \smash{$\Pic^{[0]}_{C\slash S}\to S$} is not an abelian scheme, as it is no longer proper, and we do not have an Abel--Jacobi map. There are two possible ways to fix the situations.
\begin{itemize}\itemsep=0pt
\item We can consider the compactified Jacobian \smash{$\Pic^{[0]}_{C\slash S}\subseteq \overline{\mathbb J}_{C/S}^{\mathcal E_{\pi},\sigma}$} parametrizing rank one, torsion free simple $(\mathcal E_{\pi},\sigma)$-quasi-stable sheaves on the fibers of $C\to S$ for $\mathcal E_{\pi}$ the canonical degree~$0$ polarization on $C\slash S$ in the sense of \cite[Section~2.1]{melo2015compactifications}. We refer the reader to~\cite{melo2015compactifications} and to~\cite{MeloRapagnettaViviani} for existence and properties of the compactifications.
Then a family version of the argument in \cite[Section~5]{MeloRapagnettaViviani} shows that there is an isomorphism of $S$-group
\[
\Pic^{[0]}_{C\slash S}\to \Pic^0\bigl(\overline{\mathbb J}_{C/S}^{\mathcal E_{\pi},\sigma}\bigr)
\]
and we can then define $F_*$ as the dual of $F^*$ composed with the natural inclusion of \smash{$\Pic^{[0]}_{C\slash S}$} with the compactified Jacobian.

\item Alternatively, one should in future be able to use the theory of logarithmic abelian varieties and of logarithmic Picard group developed by Molcho--Wise \cite{molchowiselog}. They show that~$\operatorname{LogPic}_{C^v\slash S}\to S$ (where the superscript $C^v$ means we are considering $C\to S$ endowed with the vertical log structure)
is a proper group stack over $\operatorname{LogSch}$
whose fibers are logarithmic abelian varieties in the sense of Kajiwara, Kato and Nakayama \cite{kajiwara2008logarithmic}.
It is expected that many classical results about abelian varieties admit generalization for logarithmic abelian varieties. In this spirit, in a forthcoming paper \cite{delignepair} Molcho--Ulirsch--Wise will show that $\operatorname{LogPic}_{C^v\slash S}$ is self dual in the category of logarithmic Abelian variety, admits a~Abel--Jacobi map and is universal for morphism to abelian groups over log schemes.
Once such a theory will be fully developed, considering the composition
 \[
 \Pic^{[0]}_{C\slash S}\hookrightarrow \operatorname{LogPic}_{C^v\slash S}\xrightarrow{F_*}\Alb(X)\times S
 \]
we would obtain the desired extension.
\end{itemize}

Once the extension $F_*$ is defined, one can look at the closed points to give an explicit description of $F_*F^*L$ and verify that this coincide with $\sum_{i=1}^n w_i a_X(f\circ q_i)$.
Since we do not need the extension at the level of universal Jacobian but only the map from $\M$, we do not fill in the technical details left out in this remark.
\end{remarkk}



\subsubsection{Refinement by correlation}\label{sec:correlationrefinement}
%\tnote{Still two moduli spaces ...}
Let us go back to the setting of $\M(X,L)$ where we fix the ramification profile. To simplify the notation, for a given point in $\M$, we set $x_i=f\circ q_i$. We are interested in the case where the contact orders have a non-trivial common divisor, i.e., there is some $1\neq\delta|\gcd(w_i)$; the we can consider the morphism
\[ a_{\bfw/\delta}\colon\ (x_i)\in X^n\longmapsto \sum_{i=1}^n \frac{w_i}{\delta}a_X(x_i). \]
Composing with the evaluation map, we define
\begin{gather}
 \kappa^\delta \colon\ \M \longrightarrow T^\delta(L,\beta)\subset\Alb(X), \nonumber\\
 \bigl(f\colon C\to X,(p_j)_{j=1}^m,(q_i)_{i=1}^n, \O_C(\alpha)\rvert_{[f]}\cong^{\lambda}f^*L\bigr) \longmapsto \sum_{i=1}^n\frac{w_i}{\delta} a_X (x_i)
.\label{eq:maptotorsion}
\end{gather}
In fact, $\kappa^\delta$ composed with the multiplication by $\delta$ is the previously defined $\kappa$. As $\kappa$ is constant equal to $\varphi_\beta(L)$ by Proposition \ref{prop-image-of-L-is-constant}, this ensures that $\kappa^\delta$ has values in $ T^\delta(L,\beta)$, the set of $\delta$-roots of $\varphi_\beta(L)$. The latter is a torsor under the finite group $\Tor_\delta(\Alb(X))$ of $\delta$-torsion elements in~$\Alb(X)$. In the particular case where $L=\mathcal O_X$, $T^\delta(X,\O_X)$ is exactly the set of $\delta$-torsion elements in the Albanese variety $\Alb(X)$.



As $L$ does not possess any canonical $\delta$-root, Proposition \ref{prop-image-of-L-is-constant} no longer applies for $\kappa^\delta$, which has no reason to be constant anymore. The morphism $\kappa^\delta$ defined in \eqref{eq:maptotorsion} thus determines a decomposition of the moduli space of logarithmic maps to $\mathbb P_X(\mathcal O_X\oplus L)$ into components, according to the value taken by $\kappa^\delta$. This refinement is called \textit{correlation}, as it stems from relations between the images of the points inside the Albanese variety. The image $\kappa^\delta(f)$ is called a \textit{correlator} and is usually denoted by $\theta$.

\begin{remarkk}
 The preimages $\bigl(\kappa^\delta\bigr)^{-1}(\theta)$ of elements $\theta\in T^\delta(L,\beta)$ are called components of the moduli space but they are not necessarily connected.
\end{remarkk}

In particular, it follows from the discussion above that we have the following.
\begin{prop}
 Let $\M$ be the moduli space of logarithmic stable maps to $Y=\PP_X(\mathcal O\oplus L)$ with the divisorial log structure defined by $D^- + D^+$. Assume that $L\in\Pic^0(X)$ and that $\delta|\gcd(w_i)$.
 Then we have a splitting of the virtual class
 \[ \vir{\M}=\sum_{\theta\in T^\delta(L,\beta)}\vir{\M^\theta},\]
 where $T^\delta(L,\beta)$ is the $\operatorname{Tor}_{\delta}(\Alb(X))$-torsor of $\delta$-roots of $\varphi_\beta(L)$.
\end{prop}

\begin{defi}\label{def:fullrefined}
We encompass the information of the correlated virtual classes in the \textit{full correlated virtual class}
\[ \fvir{\M}{\delta} = \sum_{\theta\in T^\delta(L,\beta)} \vir{\M^\theta}\cdot (\theta),\]
which is an element in the group algebra $\QQ[\Alb(X)]$ with cycle coefficients and support over the torsor $T^\delta(L,\beta)$.
\end{defi}

\subsubsection{Correlated Gromov--Witten invariants}
The projection $\PP_X(\O\oplus L)\to X$ identifies $D^-$ and $D^+$ with $X$. The evaluation map hence takes the following form:
\[
\ev\colon\ \M_{g,m}\bigl(Y|D^\pm,\beta,\bfw\bigr) \longrightarrow Y^m\times X^n,
 \]
we integrate pullback of cohomology classes over the full correlated virtual class to get \textit{correlated Gromov--Witten invariants}: for $\gamma_1,\dots,\gamma_m\in H^\bullet(Y,\QQ)$ and $\widetilde{\gamma}_1,\dots,\widetilde{\gamma}_n\in H^\bullet(X,\QQ)$, we set
\[ \bigl\langle\bigl\langle\gamma_1,\dots,\gamma_m,\widetilde{\gamma}_1,\dots,\widetilde{\gamma}_n\bigr\rangle\bigr\rangle^\delta_{g,\beta,\bfw} = \int_{\fvir{\M}{\delta}} \prod_1^m\ev_i^*(\gamma_i)\prod_1^n\widetilde{\ev}_i^*(\widetilde{\gamma}_i), \]
which is an element in $\QQ[\Alb(X)]$ with support on the torsor $T^\delta(L,\beta)$. As the computations in the elliptic case from Section \ref{sec-computation-local-invariants} illustrate, distinct correlators may provide different Gromov--Witten invariants, so that the refinement is non-trivial.






\subsection{Deformation invariants and relations of correlated classes}
Let $(\X,\L)\xrightarrow{p} B$ be a family of smooth projective varieties over $B$ and let $\mathcal L\in \Pic^0(\mathcal X\slash B)$ (in fact log smoothness is sufficient).
Let $\Y=\mathbb P(\O\oplus\L)\to B$ the associated family of $\mathbb P^1$-bundles, with the natural boundary structure $\D=\mathcal D^- +\mathcal D^-$.
Fix $\beta$ an effective curve class on $\X$ relative to~$B$ and
fix contact order $\bfw=(w_1,\dots,w_n)$ of $(q_1,\dots,q_n)$ with the boundary.
As in the previous section, the results of \cite{abramovich2014stable,gross2013logarithmic,maulik2023logarithmic} we have that the moduli stack
$\M_B:=\M_{g, m}(\mathcal Y|\D,\beta_b, \bfw)$
parametrizing diagrams of logarithmic maps
\[
\begin{tikzcd}
C\ar[r,"F"]\ar[d]\ar[d] & (\mathcal Y|\mathcal{D})\ar[d]\\
S\ar[r] & B
\end{tikzcd}
\]
is a Deligne--Mumford stack, proper over $B$ end endowed with a perfect obstruction theory $R\pi_*F^*T^{\log}_{\mathcal Y\slash B}$ relative to $B$.

As for the case of absolute stable maps,
for $B$ regular and connected, the basic properties of the virtual class construction \cite{behrend1997intrinsic,manolache2012virtual} ensure that
the degree of the class $\ev^*\gamma\cap\vir{\M_b}$ does not depend on $b\in B$; we also refer the reader to \cite[Appendix~A]{mandel2020descendant} and references therein for more details on the deformation invariance.

As $\M_B$ is a special case of Corollary~\ref{prop:logmapsdegeneration}, we still have the description of the moduli space in terms of logarithmic trivializations. As in the case of $B=\operatorname{Spec}(\mathbb C)$, we have the map
\[\kappa_B\colon\ \M_B\to\Alb(\mathcal X\slash B)\]
defined by
\[\bigl(\C\xrightarrow{f}\X,(p_j)_{j=1}^m,(q_i)_{i=1}^n,f^*\L\cong\O_{\C}(\alpha)\bigr)\rightarrow \sum_{i=1}^n w_i a_{\X\slash B}(f\circ q_i)\]
which is constant with valued $\varphi_{\beta_b}(\L_b)$ on fibers of $B$. In particular, if $\delta$ divides $\gcd(w_i)$, we get a refinement morphism
\[\kappa^\delta_B\colon\ \M_B\longrightarrow T^\delta_B(\X,\L)\subset \Alb^0(\X/B),
\]
where $T^\delta_B(\X,\L)$ is a the torsor under the $B$-group $\operatorname{Tor}_{\delta}\bigl(\Alb^0(\mathcal X\slash B)\bigr)$ of $\delta$-roots of $\varphi_{\beta_b}(\L_b)$ over~$B$. This torsor is not necessarily trivial, but we can always choose $B'\to B$ \'etale such that $T^\delta_B(\X,\L)\times_B B'\cong T^\delta_{B'}(\X',\L')$ trivializes.

Thus, working \'etale locally on $B$ we can always assume that
$T^\delta_B(\X,\L)$ has a section over~$B$, i.e., it is trivial. Then $\M_B$ split into connected component $\M_B^\theta$ indexed by the sections $\theta\colon B \to T^\delta_B(\X,\L)$.

\begin{theo}\label{thm:definvariants}
The correlated Gromov--Witten invariants are deformation invariant, i.e., let~${\X\to B}$ and $\M_B\to B$ as before and assume that $T^\delta_B(\X,\L)$ has a trivializing section $\theta$. Then~for any such section the degree of classes \smash{$\ev^*\gamma\cap\vir{\M_b^{\theta(b)}}$} for $\gamma\in\text{H}^*(\Y^{n+m})$ does not depend on~${b\in B}$.
\end{theo}

\begin{proof}
 Since \smash{$\M_B^\theta\to \M_B$} is open, the perfect obstruction theory \smash{$R\pi_*F^*T^{\log}_{\mathcal Y\slash B}$} relative to $B$ restricts to $\M_B^\theta$. The results then follows once again from the propertied of the virtual class construction \cite{behrend1997intrinsic,manolache2012virtual} as explained for example in \cite[Appendix~A]{mandel2020descendant}.
\end{proof}

By choosing suitable families of deformations of $\Y\to B$, we can exploit Theorem~\ref{thm:definvariants} to find non-trivial identities among the correlated classes.

The morphism $\varphi_\beta\colon\Pic^0(X)\to\Alb(X)$ induces a morphism between their $\delta$-torsion elements. We thus have the subgroup $\varphi_\beta\bigl(\Tor_\delta\bigl(\Pic^0(X)\bigr)\bigr)\subset\Tor_\delta(\Alb(X))$. In particular, the support $T^\delta(L,\beta)$ of the correlated Gromov--Witten invariants is stable by the action of $\varphi_\beta\bigl(\Tor_\delta\bigl(\Pic^0(X)\bigr)\bigr)$. Furthermore, we have the following.

\begin{theo}\label{theo-torsor-invariance}
 The correlated invariants \smash{$\langle\langle\gamma_1,\dots,\gamma_m,\widetilde{\gamma}_1,\dots,\widetilde{\gamma}_n\rangle\rangle^\delta_{g,\beta,\bfw}$} are invariant under the action of $\varphi_\beta\bigl(\Tor_\delta\bigl(\Pic^0(X)\bigr)\bigr)$.
\end{theo}

\begin{proof}
We consider the family over $B=\Pic^0(X)$ induced by the universal line bundle: the fiber over $L\in\Pic^0(X)$ is $\PP(\O\oplus L)$. Thus, in this case, the torsor appearing before Theorem~\ref{thm:definvariants}~is
\[ T^\delta_B(\X,\L) = \big\{ (L,\theta)\in\Pic^0(X)\times\Alb(X) \text{ s.t. } \delta\theta=\varphi_\beta(L)\in\Alb(X)\big\}.
 \]
To get a section, we use the base change $\mu_\delta\colon L\mapsto L^{\otimes\delta}$ from $\Pic^0(X)$ to itself. We thus have
\[ \mu_\delta^*T^\delta_B(\X,\L) = \big\{ (L,\theta)\in\Pic^0(X)\times\Alb(X) \text{ s.t. } \delta\theta=\varphi_\beta\bigl(L^{\otimes\delta}\bigr)\in\Alb(X)\big\}. \]
This torsor now has a section provided by $L\mapsto\theta_0(L)=\varphi_\beta(L)$.
Deformation invariance tells us that for any choice of $R$, the correlators $\theta_0(L)$ and $\theta_0(L\otimes R)$ provide the same invariants. However, the latter are correlators for the a priori distinct $\PP^1$-bundles associated to $L^{\otimes\delta}$ and~${L^{\otimes\delta}\otimes R^{\otimes\delta}}$. Assuming $R$ is of $\delta$-torsion, these $\PP^1$-bundles are the same. Furthermore, the correlators differ by
\[ \theta_0(L\otimes R)-\theta_0(L) = \varphi_\beta(L\otimes R)-\varphi(L)=\varphi_\beta(R), \]
yielding the result.
\end{proof}

Finally, we have the following relations between the classes corresponding to different levels of refinement.

\begin{lem}\label{lem-unrefinement}
The multiplication by $d$ in $\Alb(X)$ induces a morphism of $\QQ[\Alb(X)]$ denoted by~$\m{d}$. The different full refined virtual classes are related as follows:
\[ \m{\frac{\delta}{\delta'}}\bigl(\fvir{\M}{\delta}\bigr) = \fvir{\M}{\delta'}. \]
\end{lem}

\begin{proof}
 The relation merely comes from the fact that multiplication by $\delta/\delta'$ provides a surjection from $\delta$-roots of $\varphi_\beta(L)$ to $\delta'$-roots of $\varphi_\beta(L)$, so that if $\theta'\in T^{\delta'}(L,\beta)$, then we have
 \[ \M^{\theta'}=\bigsqcup_{(\delta/\delta')\theta=\theta'}\M^{\theta}. \tag*{\qed}
 \]\renewcommand{\qed}{}
\end{proof}

